%%%%%%%%%%%%%%%%%%%%%%%%%%%%%%%%%
%
%       EXERCÍCIO
%
%%%%%%%%%%%%%%%%%%%%%%%%%%%%%%%%%

\definevalorquestao

\begin{exercicioBanco}[\valorquestao]
Em um experimento, o número de bactérias presentes nas culturas A e B no instante \(t\), em horas, é dado, respectivamente, por \(A(t) = 10 \cdot 2^{t-1} + 238\) e \(B(t) = 2^{t+2} + 750\). De acordo com essas informações, o tempo decorrido, desde o início desse experimento, necessário para que o número de bactérias presentes na cultura A seja igual ao da cultura B é:

% Define as alternativas
\newcommand{\alternativas}{%
    \begin{center}
        \begin{tabularx}{\textwidth}{XXXXX}
            (a) 5 horas. &
            (b) 6 horas. &
            (c) 7 horas. &
            (d) 9 horas. &
            (e) 12 horas.
        \end{tabularx}
    \end{center}
}

% Define a resposta correta
\newcommand{\resposta}{D}

% Lógica condicional para exibição
\ifdefstring{\atividade}{lista}{%
    \alternativas
    \vspace{0.5em}
    
    \noindent\textbf{Resposta:} letra \textbf{\resposta}.
}{%
    \ifdefstring{\modo}{objetivo}{%
        \alternativas
    }{%
        % modo = discursiva → não mostra alternativas
    }
}
\end{exercicioBanco}

